\documentclass[a4paper,9pt]{extarticle}

%=========================================================
% PACKAGES
%=========================================================

\usepackage[
a4paper,
top=1.8cm,
bottom=1.5cm,
footskip=7mm,
left=1cm,
right=1cm,
headheight=14pt,
headsep=5mm
]{geometry}

\usepackage[T1]{fontenc}
\usepackage[utf8]{inputenc}
\usepackage[french]{babel}

\usepackage[sfdefault]{FiraSans}
\renewcommand{\familydefault}{\sfdefault}

\usepackage{multicol}
\setlength{\columnsep}{0.7cm}

\usepackage{amsmath,amssymb}
\usepackage{titlesec}
\usepackage{enumitem}

\usepackage{tcolorbox}
\tcbuselibrary{skins,breakable}

\usepackage{xcolor}
\usepackage{graphicx}
\usepackage{float}

\usepackage{fancyhdr}
\usepackage[hidelinks]{hyperref}


%=========================================================
% PAGE SETTINGS
%=========================================================

\setlength{\parindent}{0pt}
\setlength{\parskip}{1.5pt}

\setlength{\abovedisplayskip}{4pt}
\setlength{\belowdisplayskip}{4pt}

\setlist[itemize]{
leftmargin=*,
itemsep=1pt,
topsep=1pt
}
\setlist[enumerate]{
leftmargin=*,
itemsep=1pt,
topsep=1pt
}


%=========================================================
% COLORS
%=========================================================

\definecolor{Blue}{HTML}{E3F0FF}
\definecolor{Green}{HTML}{E3F7EA}
\definecolor{Pink}{HTML}{FFE5ED}
\definecolor{Orange}{HTML}{FFF0DC}
\definecolor{Purple}{HTML}{EEE7FF}
\definecolor{Yellow}{HTML}{FFFADB}
\definecolor{Goat}{HTML}{F7EEDF}

\definecolor{BlueTitle}{HTML}{285B8F}
\definecolor{GreenTitle}{HTML}{287052}
\definecolor{PinkTitle}{HTML}{A84063}
\definecolor{OrangeTitle}{HTML}{A96500}
\definecolor{PurpleTitle}{HTML}{6245A5}
\definecolor{GoatTitle}{HTML}{8A5A1F}
\definecolor{RecapTitle}{HTML}{5B3A8E}


%=========================================================
% HEADER
%=========================================================

\pagestyle{fancy}
\fancyhf{}

\fancyhead[L]{\small\textbf{Comportement mécanique des matériaux}}
\fancyhead[C]{\small Verres, élaboration, contraintes thermiques}
\fancyhead[R]{\small Week 05 --- 05-11/10/2026}

\fancyfoot[C]{\small\href{https://github.com/AlphaXT38G/notes_ba5}{github.com/AlphaXT38G/notes\_ba5}}
\fancyfoot[R]{\small \thepage}

\renewcommand{\headrulewidth}{0.4pt}


%=========================================================
% TITLES
%=========================================================

\titleformat{\section}
{\large\bfseries\color{BlueTitle}}
{}{0pt}{}

\titlespacing{\section}
{0pt}{6pt}{3pt}


%=========================================================
% BOXES
%=========================================================
% Boxes are kept whole (unbreakable): inside multicols, a box split by tcolorbox
% leaves an empty column behind it.

\newtcolorbox{definition}{
enhanced, unbreakable, colback=Blue, colframe=BlueTitle,
boxrule=0.7pt, arc=0mm, title=\textbf{Définition}, fonttitle=\bfseries
}

\newtcolorbox{important}{
enhanced, unbreakable, colback=Pink, colframe=PinkTitle,
boxrule=0.7pt, arc=0mm, title=\textbf{Idée clé}, fonttitle=\bfseries
}

\newtcolorbox{formula}{
enhanced, unbreakable, colback=Green, colframe=GreenTitle,
boxrule=0.7pt, arc=0mm, title=\textbf{Formules clés}, fonttitle=\bfseries
}

\newtcolorbox{warning}{
enhanced, unbreakable, colback=Orange, colframe=OrangeTitle,
boxrule=0.7pt, arc=0mm, title=\textbf{Attention}, fonttitle=\bfseries
}

\newtcolorbox{example}{
enhanced, unbreakable, colback=Purple, colframe=PurpleTitle,
boxrule=0.7pt, arc=0mm, title=\textbf{Exemple}, fonttitle=\bfseries
}

% GOAT - brown: the idea told as simply as possible
\newtcolorbox{goat}{
enhanced, unbreakable, colback=Goat, colframe=GoatTitle,
boxrule=0.7pt, arc=0mm, title=\textbf{Expliqué à une chèvre}, fonttitle=\bfseries
}

% RECAP - violet: what to know for the exam, at the end of the week
\newtcolorbox{recap}{
enhanced, breakable, colback=Purple, colframe=RecapTitle,
boxrule=1pt, arc=0mm, title=\textbf{Récap pour l'examen}, fonttitle=\bfseries
}

% Screenshot of a slide of the course: \slide{cNN-PP}
\newcommand{\slide}[2][0.95]{%
\begin{center}
\includegraphics[width=#1\linewidth]{../images/course/#2}
\end{center}}


%=========================================================
% DOCUMENT
%=========================================================

\begin{document}

\thispagestyle{fancy}

\begin{multicols}{2}

%=========================================================
\section{Les verres}

\slide[0.85]{c02-50}

\begin{definition}
Angles Si--O--Si souples $\Rightarrow$ le liquide cristallise mal ; refroidi assez vite, il se
fige : \textbf{verre}, solide \textbf{amorphe}, métastable, ordre à courte distance seulement,
ni cristaux ni joints de grains.
\end{definition}

\slide[0.8]{c02-51}

\begin{important}
$T_g$ (transition vitreuse) et densité dépendent de la vitesse de refroidissement : plus vite
$\Rightarrow$ figé plus tôt, \textbf{moins dense}.
\end{important}

\begin{formula}
\[ \tau = \eta\dot\gamma = \eta\frac{d\gamma}{dt}, \qquad \sigma = 3\eta\dot\varepsilon, \qquad
   \eta = K\,e^{+Q/RT} \]
Eau $10^{-3}$, huile $5\cdot10^{-2}$, miel 2--10\,Pa$\cdot$s. Repères du verre : fusion $10$,
mise en forme $10^3$, ramollissement $4\cdot10^6$\,Pa$\cdot$s.
\end{formula}

\begin{important}
\textbf{Modifiants} (Na$_2$O, CaO, B$_2$O$_3$, quelques dizaines de \%) : cassent le réseau,
baissent la viscosité, la silice devient formable plus bas.
\end{important}

\begin{goat}
Un verre est un liquide figé en désordre. Chaud, il coule comme du miel, de plus en plus
facilement quand on chauffe ; la soude coupe des liens pour qu'il coule plus tôt.
\end{goat}

%=========================================================
\section{Élaboration}

\begin{itemize}
\item \textbf{Argiles} : formage humide (malaxeur, extrusion, calibrage), \textbf{barbotine}
      (moule en plâtre qui boit l'eau), \textbf{pressage} de poudres ; séchage (retrait puis
      arrêt) ; cuisson.
\item \textbf{Capillarité} : énergie de surface $\gamma$ (J/m$^2$ = N/m).
\item \textbf{Frittage} : à chaud, la poudre réduit sa surface ; cous, pores qui rétrécissent et
      disparaissent ; densification sans fusion.
\item Ciment, plâtre : durcissement \textbf{chimique}.
\item \textbf{Verre} : newtonien $\Rightarrow$ pas de striction $\Rightarrow$ étirage (fibres),
      pressage--soufflage (bouteilles), float (bain d'étain).
\end{itemize}

\slide[0.8]{c02-67}

\begin{goat}
Une poudre a trop de surface ; en chauffant, les grains se collent et les trous se referment.
Le verre chaud s'étire en fil sans jamais s'amincir à un endroit.
\end{goat}

%=========================================================
\section{Contraintes thermiques}

\begin{formula}
\[ \sigma = \frac{E}{1-\nu}\,\alpha\,(T_1 - T_2) \]
Couche de surface refroidie d'un coup, collée à un c\oe ur qui ne se déforme pas ;
\textbf{traction} ; $1-\nu$ car état \textbf{biaxial}.
\end{formula}

\begin{example}
Verre sodocalcique : $E = 74$\,GPa, $\nu = 0{,}25$, $\alpha = 8{,}5\cdot10^{-6}$, $\Delta T = 200$\,K
$\Rightarrow \sigma \approx 170$\,MPa : \textbf{choc thermique}, il casse.
\end{example}

\begin{warning}
La \og résistance au choc thermique\fg{} ($\Delta T$ max à la trempe) n'est pas une propriété du
matériau : défauts, forme, taille, vitesse comptent ; procédure normalisée.
\end{warning}

\slide[0.8]{c02-80}

\begin{important}
\textbf{Verre trempé} : refroidi encore mou, le c\oe ur suit la surface, puis se contracte en
dernier : \textbf{surface en compression} (protège des fissures), c\oe ur en traction
(émiettement si une fissure l'atteint). \textbf{Échange d'ions} (K$^+$ pour Na$^+$) : $> 500$\,MPa
sur une fine couche (Gorilla).
\end{important}

\begin{goat}
Peau refroidie d'un coup : elle veut rétrécir, le c\oe ur la retient, elle se fissure. Si le
c\oe ur est encore mou, c'est lui qui rétrécit en dernier et serre la peau : une peau serrée ne se
fissure pas.
\end{goat}

%=========================================================
\section{Récap de la semaine}

\begin{recap}
\begin{itemize}
\item Verre : amorphe, métastable ; $T_g$ et densité selon la vitesse.
\item $\tau = \eta\dot\gamma$, $\sigma = 3\eta\dot\varepsilon$ ; $\eta = Ke^{Q/RT}$ ;
      $10$ / $10^3$ / $4\cdot10^6$\,Pa$\cdot$s ; modifiants.
\item Argiles, barbotine, pressage, séchage ; $\gamma$ ; frittage ; ciment chimique ; verre :
      pas de striction, étirage, soufflage, float.
\item $\sigma = \frac{E}{1-\nu}\alpha(T_1 - T_2)$ (170\,MPa) ; choc thermique ; verre trempé,
      échange d'ions.
\end{itemize}
\end{recap}

\end{multicols}

\end{document}
